\clearpage
\fancyhead[L]{\small CHRONOMETRISCHE LÄNGENBESTIMMUNGEN.}
\hypertarget{native-460-chronometric}{}
\expandafter\def\csname HistoricalNoteText462\endcsname{[beobachtet von Dr. \textsc{Tiarks} auf der von der englischen Admiralität im Jahre 1824
veranstalteten Expedition zur chronometrischen Bestimmung der Längsunterschiede zwischen Altona,
Bremen, Helgoland und Greenwich. \textsc{Schering}.]}
\begin{center}{\large CHRONOMETRISCHE LÄNGENBESTIMMUNGEN.}\end{center}

\begin{center}
\rule{0.55\linewidth}{0.4pt}\\[2pt]
Astronomische Nachrichten. Band V. Nr. 110. Seite 227\,..\,234. 1826 November.\\[2pt]
\rule{0.55\linewidth}{0.4pt}
\end{center}

Es seien $\theta, \theta', \theta''\ldots$ die Zeiten (insgesamt $n$), zu denen der Chronometer der jeweiligen Ortszeit an den Orten mit den Längen $x, x', x''$ usw. um die Unterschiede $a, a', a''$ usw. voraus war. Ich setze voraus, dass die Angaben $\theta, \theta', \theta''$ usw. bereits auf einen einzigen Ort reduziert sind. Ist also der tägliche Gang des Chronometers $=u$, so hätte man bei einem vollkommenen Chronometer die $n-1$ Gleichungen.

\[
a-\theta u-x = a'-\theta' u-x' = a''-\theta'' u-x'' = a'''-\theta''' u-x''' = \text{etc.}
\]

\begingroup\rightskip=0pt plus 1em
Damit diese Gleichungen ausreichen, um die unbekannten Größen $u, x, x', x''$ usw. zu bestimmen, betrachtet man eine der Größen $x, x', x''$ usw. als gegeben und setzt voraus, dass mindestens an einem Ort zweimal beobachtet wurde, sodass zwei der Größen $x, x', x''$ usw. identisch sind. Sind nicht mehr als zwei identisch, ist die Aufgabe vollständig bestimmt. Im anderen Fall ist sie überbestimmt. Dann muss man die unbekannten Größen so bestimmen, dass die $n-1$ Gleichungen\par\endgroup

\[
\begin{aligned}
0 &= a-a'+(\theta'-\theta)u-x+x'\\
0 &= a'-a''+(\theta''-\theta')u-x'+x''\\
0 &= a''-a'''+(\theta'''-\theta'')u-x''+x''' \quad \text{etc.}
\end{aligned}
\]
\noindent so genau wie möglich erfüllt sind. Denn wegen der Unvollkommenheiten, die bei allen Chronometern stets vorkommen, wird man nicht jede Gleichung genau erfüllen können. Diesen Gleichungen darf man aber offensichtlich nicht dasselbe Gewicht geben. Die Größen

\[
\begin{aligned}
&a-a'+(\theta'-\theta)u-x+x'\\
&a'-a''+(\theta''-\theta')u-x'+x''
\end{aligned}
\]
\noindent drücken nämlich nur die Summen aller Abweichungen vom mittleren Gang aus, die der Chronometer in den Zeitintervallen $\theta'-\theta$, $\theta''-\theta'$ usw. hatte. Bei einem guten Chronometer, dem man tatsächlich einen mittleren Gang zuschreiben kann, der keinen allmählich in derselben Richtung zunehmenden Änderungen unterliegt, muss man den zu erwartenden mittleren Wert einer solchen Summe proportional zur Quadratwurzel des Zeitintervalls ansetzen.

\Needspace{3\baselineskip}
\SourceLeaf{461}
Daher muss man bei der Behandlung der obigen Gleichungen nach der Methode der kleinsten Quadrate ungleiche Gewichte ansetzen, die zu den Zeitintervallen $\theta'-\theta$, $\theta''-\theta'$, $\theta'''-\theta''$ usw. umgekehrt proportional sind.

Die Auflösung bereitet dann keine Schwierigkeiten, und man erhält sowohl die plausibelsten Werte von $u, x, x', x''$ usw. als auch ihre relative Zuverlässigkeit. Dazu mache ich noch einige Bemerkungen.

1) Sind die erste und die letzte Beobachtung am selben Ort gemacht worden, so stimmt der plausibelste Wert von $u$ genau mit dem Wert überein, der allein aus dem Vergleich der beiden äußersten Beobachtungen folgt. Die Rechnung wird dann außerordentlich einfach. Nach einem leicht zu beweisenden Lehrsatz darf man diesen plausibelsten Wert von $u$ unmittelbar in die Gleichungen einsetzen. Das ist gleichbedeutend damit, sämtliche beobachteten Chronometerzeiten auf die Zeiten eines gedachten Chronometers zu reduzieren, dessen Vorlauf $=0$ wäre.

2) Hat man den Gleichungen einfach die Gewichte $\dfrac{1}{\theta'-\theta}$, $\dfrac{1}{\theta''-\theta'}$ gegeben, so ist die Einheit der Gewichte, die man für die Endergebnisse der Längenbestimmungen findet, die Genauigkeit, die man mit \emph{diesem} Chronometer bei bekanntem Gang für einen nach einem Zeitintervall von einem Tag bestimmten Längenunterschied erwarten würde (sofern die Zeiten $\theta, \theta', \theta''$ in Tagen ausgedrückt sind). Um jedoch die Ergebnisse \emph{verschiedener} Chronometer von ungleicher Güte vergleichen zu können, braucht man noch einen Faktor, der von der Güte des jeweiligen Chronometers abhängt. Um ihn zu finden, setzt man die Werte der Größen

\[
\begin{aligned}
&a-a'+(\theta'-\theta)u-x+x'\\
&a'-a''+(\theta''-\theta')u-x'+x''\\
&a''-a'''+(\theta'''-\theta'')u-x''+x''' \quad \text{etc.}
\end{aligned}
\]
\noindent nachdem man für $u, x, x', x''$ usw. die gefundenen plausibelsten Werte eingesetzt hat, $=\lambda, \lambda', \lambda''$ usw. und

\[
\frac{\lambda\lambda}{\theta'-\theta} + \frac{\lambda'\lambda'}{\theta''-\theta'} + \frac{\lambda''\lambda''}{\theta'''-\theta''} + \text{etc.} = S
\]

Ferner sei $\nu$ die Anzahl sämtlicher zuvor unbekannter Größen und $m=\BareRoot{\dfrac{S'}{n-\nu-1}}$. Dann ist der für den jeweiligen Chronometer spezifische Faktor der Größe $\dfrac{1}{mm}$ oder $\dfrac{n-\nu-1}{S}$ proportional. Man kann $m$ als die nach einem Zeitintervall von einem Tag zu erwartende mittlere Abweichung vom mittleren Gang ansehen.

3) Die obigen Regeln gelten für einen Chronometer, dessen Gang sich nicht erheblich fortschreitend ändert. Im gegenteiligen Fall kann man sich, sofern die Beobachtungsreihe nicht übermäßig lang ist, damit begnügen, eine zur Zeit proportionale Änderung des täglichen Ganges anzunehmen. Dann ist eine weitere unbekannte Größe einzuführen, und die Gleichungen haben folgende Gestalt:

\[
\begin{aligned}
0 &= a-a'+(\theta'-\theta)u+(\theta'\theta'-\theta\theta)v-x+x'\\
0 &= a'-a''+(\theta''-\theta')u+(\theta''\theta''-\theta'\theta')v-x'+x''\\
0 &= a''-a'''+(\theta'''-\theta'')u+(\theta'''\theta'''-\theta''\theta'')v-x''+x'''\\
&\text{usw.}
\end{aligned}
\]

4) Noch weiter zu gehen und eine weitere unbekannte Größe sowie Glieder der Form $(\theta'^3-\theta^3)w$ einzuführen, dürfte kaum ratsam sein. Unter anderem würde ich Chronometer, die starke, deutliche, aber ihrerseits unregelmäßige Änderungen des mittleren Ganges zeigen, lieber ganz ausschließen. Ihre Ergebnisse sind nämlich wesentlich ungenauer, und ihre \SourceLeaf{462}Genauigkeit lässt sich zum Vergleich viel schwerer in Zahlen angeben. Ich gehe daher hier auf die viel verwickeltere Theorie solcher Fälle nicht ein, da im vorliegenden Fall das Obige ausreichen wird.

4) Zur Auflösung der Gleichungen nach der Methode der kleinsten Quadrate ist es vielleicht nicht überflüssig, daran zu erinnern, dass man in den meisten Fällen die unbekannten Größen zweckmäßig aus einem bekannten (möglichst gut angenäherten) und einem unbekannten (also sehr kleinen) Teil zusammensetzt. Dieser Rat wurde zwar schon mehrfach an anderer Stelle gegeben, und der Vorteil dieses Vorgehens liegt auf der Hand. Dennoch schien es mir gut, daran zu erinnern, weil ich sehe, dass dieser Rat häufig vergessen wird. Dadurch werden die numerischen Rechnungen unnötig erschwert, und Fehler können leichter auftreten.

\Needspace{23\baselineskip}
Von den 36 Chronometern\HistoricalNote{462} habe ich die Ergebnisse für die folgenden 5 berechnet.

\begin{center}
\begingroup\fontsize{8}{10}\selectfont\setlength{\tabcolsep}{2pt}
\begin{tabular}{ll|r|r|r|r|r|r}
 & & & Nr. 1 & Nr. 4 & \textsc{Breguet} 3056 & \textsc{Kessels} 1252 & \textsc{Barraud} 904 \\
\hline
Greenwich & Juni 30 & $3^{\rm u}22^{\rm m}$ & $-8^{\rm m}17^{\rm s}14$ & $+1^{\rm m}2^{\rm s}37$ &  &  &  \\
 & Juli 25 & $2\;15$ & $10\;44.39$ & $1\;32.15$ & $+30^{\rm m}59^{\rm s}75$ & $+50^{\rm m}29^{\rm s}31$ & $+48^{\rm m}29^{\rm s}20$ \\
 & 28 & $3\;13$ & $11\;0.69$ & $1\;36.96$ & $30\;50.07$ & $50\;39.69$ & $48\;40.24$ \\
 & Aug. 2 & $1\;15$ & $11\;28.48$ & $1\;44.44$ & $30\;31.78$ & $50\;52.14$ & $48\;58.87$ \\
 & 17 & $10\;28$ & $12\;59.40$ & $2\;6.24$ & $29\;35.69$ & $51\;38.66$ & $49\;57.83$ \\
 & 25 & $7\;27$ & $13\;47.98$ & $2\;15.84$ & $29\;10.48$ & $52\;2.45$ & $50\;27.15$ \\
 & Sept. 10 & $7\;40$ & $15\;24.47$ & $2\;40.36$ &  &  &  \\
Helgoland & Juli 3 & $3\;40$ & $-40\;8.00$ & $-30\;26.84$ &  &  &  \\
 & 22 & $12\;40$ & $42\;2.02$ & $30\;3.89$ & $-0\;20.34$ & $+18\;48.39$ & $+16\;47.39$ \\
 & Aug. 5 & $1\;48$ & $43\;18.11$ & $29\;43.35$ & $1\;10.24$ & $19\;26.77$ & $17\;37.51$ \\
 & 11 & $13\;9$ & $43\;35.77$ & $29\;33.43$ & $1\;32.75$ & $19\;47.22$ & $18\;1.30$ \\
 & 30 & $19\;30$ & $45\;53.08$ & $29\;7.96$ & $2\;40.67$ & $20\;47.68$ & $19\;17.03$ \\
 & Sept. 6 & $3\;6$ & $46\;31.56$ & $28\;58.94$ & $3\;4.55$ & $21\;6.56$ & $19\;43.80$ \\
 & 7 & $8\;42$ & $46\;38.72$ & $28\;56.71$ &  &  &  \\
Altona & Aug. 6 & $5\;55$ & $-51\;38.95$ & $-37\;55.76$ & $-9\;28.50$ & $+11\;16.25$ & $+9\;28.48$ \\
 & 9 & $12\;35$ & $51\;57.35$ & $37\;50.03$ & $9\;38.81$ & $11\;27.76$ & $9\;40.30$ \\
 & 31 & $9\;57$ & $54\;10.33$ & $37\;21.30$ & $10\;56.68$ & $12\;35.96$ & $11\;5.92$ \\
 & Sept. 4 & $22\;12$ & $54\;39.16$ & $37\;15.21$ & $11\;15.36$ & $12\;48.10$ & $11\;24.49$ \\
Bremen & Aug. 13 & $0\;2$ & $-47\;50.65$ & $-33\;16.49$ & $-5\;23.37$ & $+16\;5.83$ & $+14\;21.86$ \\
\end{tabular}
\endgroup
\end{center}

Zur Probe gebe ich die Rechnung für \emph{Breguet} 3056 an. Die Länge von Helgoland sei $=0$, die von Greenwich $=-x$ und die von Altona $=+y$. Bremen schließe ich hier aus, weil dort nur einmal beobachtet wurde und es daher ohnehin keine Kontrolle ermöglicht.

Die Zeiten rechne ich vom ersten Vergleich der englischen Chronometer an (Greenwich, 30. Juni, $3^{\rm u}22^{\rm m}$).

Indem ich einen gedachten Chronometer mit dem Gang $=0$ einsetze, finde ich seinen Stand.

\begin{center}
\begin{tabular}{r|l||r|l}
$\theta$ & & $\theta$ & \\
\hline
22.4 & $+60^{\rm s}20$ & 42.4 & $+59^{\rm s}88$ \\
25.0 & $+1949.60-x$ & 48.3 & $+1949.60-x$ \\
28.0 & $+1950.87-x$ & 56.2 & $+1952.74-x$ \\
32.9 & $+1950.29-x$ & 61.6 & $+61.32$ \\
35.9 & $+59.08$ & 62.2 & $-432.53+y$ \\
37.1 & $-434.98+y$ & 66.8 & $-434.98+y$ \\
40.4 & $-433.49+y$ & 68.0 & $+60.19$ \\
\end{tabular}
\end{center}

\SourceLeaf{463}
\Needspace{12\baselineskip}
Die obigen Gleichungen sind hier so beschaffen, dass $x$ und $y$ überhaupt nicht vermischt sind. Dadurch wird die weitere Rechnung noch einfacher. Für $x$ haben wir nämlich vier Bestimmungen.

\[
\begin{array}{ll}
+1889^{\rm s}40 & \text{Gewicht } \dfrac{1}{2.6}=0.38 \\[7pt]
+1891.21 & \dfrac{1}{3.0}=0.33 \\[7pt]
+1889.78 & \dfrac{1}{5.9}=0.17 \\[7pt]
+1891.42 & \dfrac{1}{3.4}=0.19 \\[7pt]
\hline
\text{Also } x=+1890^{\rm s}36 \quad \text{Gewicht} & =1.07
\end{array}
\]

Ebenso findet man

\[
y=+494^{\rm s}12 \quad \text{Gewicht } =3.83
\]
\noindent Setzt man diese Werte ein, so erhält man den Stand des gedachten Chronometers gegenüber der Helgoländer Zeit

\begin{center}
\begingroup\small\renewcommand{\arraystretch}{1.25}
\begin{tabular}{r|r r@{\qquad}r|r r}
$\theta$ & & $\lambda$ & $\theta$ & & $\lambda$ \\
\cline{1-1}\cline{3-4}\cline{6-6}
$22.4$ & $+60^{\rm s}20$ & \raisebox{-7pt}{$-0^{\rm s}96$} & $42.4$ & $+59^{\rm s}88$ & \raisebox{-7pt}{$-0^{\rm s}62$} \\
$25.0$ & $59.24$ & \raisebox{-7pt}{$+1.27$} & $48.3$ & $59.24$ & \raisebox{-7pt}{$+3.14$} \\
$28.0$ & $60.51$ & \raisebox{-7pt}{$-0.58$} & $56.2$ & $62.38$ & \raisebox{-7pt}{$-1.06$} \\
$32.9$ & $59.93$ & \raisebox{-7pt}{$-0.85$} & $61.6$ & $61.32$ & \raisebox{-7pt}{$+0.27$} \\
$35.9$ & $59.08$ & \raisebox{-7pt}{$+0.06$} & $62.2$ & $61.59$ & \raisebox{-7pt}{$-2.45$} \\
$37.1$ & $59.14$ & \raisebox{-7pt}{$+1.47$} & $66.8$ & $59.14$ & \raisebox{-7pt}{$+1.05$} \\
$40.4$ & $60.63$ & \raisebox{-7pt}{$-0.75$} & $68.0$ & $60.19$ &  \\
$42.4$ & $59.88$ &  &  &  &  \\
\end{tabular}
\endgroup
\end{center}

Also $S=6.00$, $m=\BareRoot{\dfrac{6.00}{13-3}}$. Daraus ergibt sich der zu erwartende mittlere Fehler für $x$: $0^{\rm s}75$, für $y$: $0^{\rm s}40$

\Needspace{19\baselineskip}
Die 5 Chronometer, für die ich die Rechnung durchgeführt habe, ergeben insgesamt

\begin{center}
\begin{tabular}{ll|l|l}
 & & E.\,med. & Gewicht \\
\hline
Breguet & $x=1890^{\rm s}36$ & $0^{\rm s}75$ & 1.78 \\
Kessels & $1893.29$ & $0.67$ & 2.23 \\
Barraud & $1892.32$ & $0.49$ & 4.16 \\
Engl. 1 & $1892.39$ & $0.43$ & 5.41 \\
---\quad 4 & $1892.52$ & $0.35$ & 8.16 \\
\hline
Mittel & $x=1892.35$ & & 21.74 \\
\end{tabular}
\end{center}

\begin{center}
\begin{tabular}{ll|l|l}
Breguet & $y=494.12$ & $0.40$ & 6.25 \\
Kessels & $493.89$ & $0.36$ & 7.72 \\
Barraud & $493.67$ & $0.26$ & 14.79 \\
Engl. 1 & $493.98$ & $0.29$ & 11.89 \\
---\quad 4 & $494.16$ & $0.24$ & 17.36 \\
\hline
 & $y=493.96$ & & 58.01 \\
\end{tabular}
\end{center}

\begingroup\emergencystretch=1.5em\relax
In die letzte Spalte mit der Überschrift Gewicht habe ich\newline hier $\dfrac{1}{\text{Quadr. (E.\,m.)}}$ eingetragen und somit als Einheit die Genauigkeit angesetzt, für die der zu erwartende mittlere Fehler $=1^{\rm s}$ ist. So ergibt sich zum Beispiel für Altona als zu erwartender mittlerer Fehler $=\dfrac{1^{\rm s}}{\BareRoot{58.01}}=0^{\rm s}13$. Es ist allerdings ratsamer, die Zahlen der letzten Spalte nur als Verhältniszahlen zu betrachten und die absolute Genauigkeit aus den Abweichungen der mit den einzelnen Chronometern für $x$ und $y$ gefundenen Werte von den Endergebnissen abzuleiten. Doch wird die Genauigkeit des Endergebnisses auch dann \SourceLeaf{464}noch etwas größer erscheinen, als sie wirklich ist, da die Zeitbestimmungen in Greenwich, Helgoland und Altona keine \emph{absolute} Genauigkeit haben. Deshalb müssen sich die aus dieser Quelle stammenden Fehler offenbar selbst bei noch so vielen Chronometern in den Endergebnissen auswirken.
\par\endgroup

Die Länge von Bremen lässt sich auf folgende Weise bestimmen. Setzt man die Länge $=z$ östlich von Helgoland an, so ergibt der Vergleich des Chronometers von \textsc{Breguet}

\[
-165^{\rm s}52+z
\]


\begin{center}\small
\begin{tabular}{p{0.42\linewidth} r l r}
Also aus dem vorhergehenden Vergleich & $z=225.40$ & Gewicht & $\dfrac{1}{1.4}=0.7$ \\[5pt]
aus dem folgenden Vergleich & $z=224.76$ & & $\dfrac{1}{4.5}=0.2$ \\[8pt]
\cline{2-2}\cline{4-4}
 & $225.24$ & & $0.9$ \\
\end{tabular}
\end{center}

\Needspace{17\baselineskip}
Das Gewicht $0.9$ muss noch mit $\dfrac{10}{6.00}$ multipliziert werden. So ergeben die 5 Chronometer

\begin{center}
\begin{tabular}{l|l|l}
 & & Gewicht \\
\hline
Breguet & $225.24$ & $1.5$ \\
Kessels & $225.84$ & $1.9$ \\
Barraud & $225.39$ & $3.6$ \\
Engl. 1 & $226.04$ & $2.9$ \\
---\quad 4 & $224.86$ & $4.3$ \\
\hline
 & $225.42$ & $14.2$ \\
\end{tabular}
\end{center}

Die Länge von Bremen, die hiernach $268^{\rm s}54$ westlich von Altona liegt, hängt aber natürlich weiterhin von der Zeitbestimmung in Bremen ab, und dieser Unterschied scheint mehrere Sekunden zu klein zu sein. Nach meinen Dreiecken liegt der Ansgariusthurm $273^{\rm s}51$ in Zeit westlich von Göttingen, also das Observatorium von \textsc{Olbers} $271^{\rm s}9$.
